\section{Discussion}\label{sec:discussion}\justifying{}
In this study, a physiologically based pharmacokinetic (PBPK) model for enalapril was developed and applied under various physiological and pathological conditions. Using a comprehensive database derived from 49 clinical trials, this model reliably predicts the pharmacokinetics of enalapril and its active metabolite, enalaprilat, in a diverse patient population. The focus on renal and hepatic impairment is particularly important given their impact on drug clearance and metabolism, which can significantly affect therapeutic outcomes.

% Database}
This work has created the first openly available comprehensive database of enalapril pharmacokinetics and pharmacodynamics. To our knowledge, this is the first time that systematic data curation of enalapril pharmacokinetic data has been undertaken. All data are freely available for re-use and can be downloaded from \href{https://pk-db.com}{https://pk-db.com}.

% PBPK Model
Despite the developed PBPK model being in good agreement with the experimental data certain limitations exist. (i) The focus of the model was on the description of the pharmacokinetics of enalapril and enalaprilat. Pharmacodynamic data was curated on systolic and diastolic blood pressure, heart rate, mean arterial pressure and substances such as renin, angiotensin I and II and aldosterone from the 49 studies, but is not part of the computational model. (ii) The model of renal and hepatic functional impairment was based on existing validated models for indocyanine green~\cite{Koller2021, Koller2021a} and talinolol~\cite{Mallol2023}. In contrast the cardiac functional impairment model was very simplified and only assumed reduced cardiac output. A more detailed model of the pathophysiological changes in cardiac impairment could substantially improve the agreement between predictions and clinical data. (iii) The tissue models for renal excretion, hepatic metabolism and intestinal uptake are simplifications, mainly due to the limited availability of data. I.e. transport and metabolism were modeled as simple mass-action or Michaelis-Menten kinetics often lumping multiple steps in a single overall reaction rate. With the availability of more detailed experimental information these models could be extended.

% Parameter optimization
In this work, a multi-step approach was used to optimise the parameters. I.e. first the parameters specific for enalaprilat were optimised, then the parameters for enalapril and finally the parameters for the hepatic conversion of enalapril to enalaprilat. All three optimisations significantly improved the agreement between the model and the experimental data, and the resulting model is in good agreement with most of the training data. I.e. the optimal parameters given in Tab.~\ref{tab:Parameter_Fitting} provide a good model fit as shown in Fig.~\ref{fig:fitting_result_ena}, Fig.~\ref{fig:fitting_result_eat} and Fig.~\ref{fig:fitting_result_enaeat}. 

There are some limitations in the parameter optimisation: (i) The optimisation used a local gradient descent combined with a multistart method (n=100). We did not use a global optimisation approach such as differential evolution, which could give better results. Most of the optimisation runs resulted in identical optimal parameters (as seen in the waterfall plots), with some runs resulting in two alternative local optima. From our experience with previous models, the local multi-start methods with a large number of runs and the single global optimisations gave very similar results, although the global optimisations were much more computationally intensive. Therefore, we restricted the parameter optimisation to a local optimiser. (ii) A multi-step method was used instead of a single-step method that fits all parameters with all training data at once. The main reason for this is the heterogeneity of the data. Most studies report plasma or serum time courses of enalaprilat, enalapril is reported much less frequently, and data on urinary levels of enalapril and enalaprilat are much more limited. As a result of this imbalance in exercise data, certain parameters are difficult to determine. It is necessary to weight the data to account for the imbalance in the exercise data. Instead of implementing a weighting scheme, we decided to perform a multi-step optimisation, which results in much less bias in the training data, rather than pooling all the data in a single step method. A single global optimisation of all parameters at once could provide a better overall fit because the parameters in the multi-step methods could compensate for each other. (iv) Classical single point optimisation was used instead of the Bayesian approach. This could provide a better fit and an estimate of the variability (posteriors) of the parameters, but at the cost of a much higher computational cost. 

Despite these limitations, parameter optimisation yielded strong agreement between model predictions and experimental data. Discrepancies between simulation and data remained for some studies, notably Dickstein et al. 1987~\cite{Dickstein1987} and Hockings et al. 1986~\cite{Hockings1986}. While the model could plausibly be refined to better reproduce these data points, the reported data themselves are in some cases unreliable: assays often have limited accuracy at low enalapril and enalaprilat concentrations, are occasionally insufficiently specific to distinguish enalapril from enalaprilat and can give false positives due to cross-reactivity, and are subject to systematic reporting errors. Even where the simulation and reference data disagree, the model therefore still provides a robust estimate of the underlying pharmacokinetics. For a discussion and examples of these problems see Grzegorzewski et al.~\cite{Grzegorzewski2021a}.


% Simulations
The developed PBPK model accurately simulates a wide range of clinical studies of enalapril in different patient populations, single and multiple dose, oral and intravenous routes, and different pathophysiologies (renal and hepatic impairment). As shown in Fig.~\ref{fig:DoseDependency}, Fig.~\ref{fig:Schwartz1985_HT} and Fig.~\ref{fig:Schwartz1985_CHF}, the model could adapt to a range of dosing variations. It could also adapt to renal and hepatic impairment and give a good estimate of the effect of these two conditions on the pharmacokinetics of enalapril. 

The model was able to capture the reduced elimination of enalapril in patients with renal impairment with a good degree of accuracy. The simulation could not match the observed reference data perfectly, but it captured the overall trend very well. (Fig.~\ref{fig:Kelly1986_single}, Fig.~\ref{fig:Kelly1986_multi}, Fig.~\ref{fig:Fruncillo1987}). Importantly, these simulations are model validations and none of the disease data were used for parameter optimisation. Including some of the renal impairment data in the model fitting would improve the fit.

In the case of hepatic impairment, the model was able to account for the reduced metabolic capacity of the liver and simulate the pharmacokinetics of enalapril and enalaprilat in plasma and urine with very good agreement with reference data (Fig.~\ref{fig:Baba1990}, Fig.~\ref{fig:Ohnishi1989}).

Mutations in the CES1 genotype are an important factor affecting the pharmacokinetics of enalapril and the model was able to reproduce the effects through simulations (Fig.~\ref{fig:Her2021_ces1}, Fig.~\ref{fig:Stage2017_ces1_GAIN}, Fig.~\ref{fig:Stage2017_ces1_LOSS}, Fig.~\ref{fig:Tarkiainen2015_ces1}). The overall effect of reduced and increased activity of CES1 due to variants was captured by the model. Due to the large variability between different studies, the predicted curves have a systematic bias. Importantly, none of the CES1 variant data were used in parameter optimisation, i.e. all predictions are model checks. With appropriate parameterisation, the model will be able to simulate the known variants even better, and in the future it could simulate additional CES1 variants as activity data of the variants become available. 

% Parameter scans
The predicted changes in enalapril and enalaprilat pharmacokinetics with changes in renal function, hepatic function and CES1 activity are in good agreement with cinical data.

Plasma enalapril concentrations decreased with increasing renal activity and decreasing hepatic impairment (Fig.~\ref{fig:timecourse_hepatic_renal}), reflecting more efficient renal excretion and greater hepatic conversion of enalapril to enalaprilat, respectively. The same pattern was evident in urine, where enalaprilat concentrations were higher in patients with less severe cirrhosis. As hepatic impairment increased, enalapril concentrations rose while enalaprilat concentrations fell, confirming that the model correctly captures hepatic metabolism of enalapril to enalaprilat. With increasing renal impairment, both enalapril and enalaprilat concentrations increased, consistent with impaired renal elimination of both compounds. CES1 activity produced the expected opposing trends: higher CES1 activity accelerated both enalapril elimination and enalaprilat formation, whereas reduced CES1 activity had the opposite effect, as captured by the parameter scans.

Our analysis focused on the effect of renal function, liver function and CES1 activity. No systematic parameter sensitivity analysis was performed, but the focus was on the effect of selected parameters corresponding to typical functional changes, i.e. hepatic impairment, renal impairment, changes in CES1 activity. A more detailed analysis of the effect of different parameters on the pharmacokinetics of enalapril may provide additional information.

% Conclusion
In conclusion, this study has developed and validated a physiologically based pharmacokinetic (PBPK) model for enalapril, which has significantly improved our understanding of its pharmacokinetics under various physiological and pathological conditions. Using data from 49 clinical trials, the model provides robust predictions of the behaviour of enalapril and enalaprilat in a wide range of patient populations. Its accuracy in simulating plasma and urine pharmacokinetic parameters in healthy individuals, as well as its ability to reflect changes due to renal and hepatic impairment and CES1 genetic variants, demonstrates its practical relevance and scientific validity. The predictive performance of this model is not only consistent with the existing literature, but also provides a reliable framework for further research and application in clinical settings. Thus, this work represents a substantial contribution to pharmacokinetic modelling, paving the way for more personalised and effective therapeutic strategies in the treatment of cardiovascular disease and beyond.